\documentclass[11pt]{article}
\usepackage[a4paper,margin=18mm]{geometry}
\usepackage{fontspec}
\setmainfont{Latin Modern Roman}
\usepackage{amsmath,amssymb,amsthm,amscd,graphicx,color}
\usepackage{xurl}
\usepackage[unicode,hidelinks]{hyperref}
\allowdisplaybreaks
\setlength\emergencystretch{3em}
\numberwithin{equation}{section}

\newtheorem{Theorem}{Theorem}[section]
\newtheorem{Corollary}[Theorem]{Corollary}
\newtheorem{Lemma}[Theorem]{Lemma}
\newtheorem{Proposition}[Theorem]{Proposition}
 { \theoremstyle{definition}
\newtheorem{Definition}[Theorem]{Definition}
\newtheorem{notation}[Theorem]{Notation}
\newtheorem{Example}[Theorem]{Example}
\newtheorem{Remark}[Theorem]{Remark} }


\DeclareMathOperator{\rig}{rig}
\DeclareMathOperator{\Aut}{Aut}
\DeclareMathOperator{\dR}{dR}
\DeclareMathOperator{\MW}{MW}
\DeclareMathOperator{\diag}{diag}
\DeclareMathOperator{\prim}{prim}
\DeclareMathOperator{\Frob}{Frob}
\DeclareMathOperator{\sig}{sign}



\begin{document}
\begin{center}\Large Common Frobenius characteristic factor for tame Delsarte monomial deformations with a common cover and injective complement restriction\end{center}
\noindent\textbf{Stable ID:} 1888\par
\noindent\textbf{Authors:} Remke Kloosterman\par
\noindent\textbf{Original work:} Zeta Functions of Monomial Deformations of Delsarte Hypersurfaces\par
\noindent\textbf{Version:} Official SIGMA 13 (2017) article 087, DOI 10.3842/SIGMA.2017.087; journal PDF and native TeX acquired 2026-10-01, exact bytes pinned by SHA256\par
\noindent\textbf{Selected task scope:} Selected characteristic-polynomial conclusion (last sentence) of Theorem2.18; the additional module quotient-direction assertion is outside the selected scope. Theorem2.18: n>=\allowbreak{}2 over finite F\_q; each coefficient matrix A\_i invertible with nonnegative integral entries, a zero in each column and positive A\_i\textasciicircum{}-1(1,...,1)\textasciicircum{}T. Deformation vectors a\_i are nonnegative integral of weighted degree d and a\_i A\_i\textasciicircum{}-1 nonnegative; gcd(q,d)=\allowbreak{}1 for a common integral denominator d, and gcd(q,(n+\allowbreak{}1)det(A\_i))=\allowbreak{}1 as additionally required by source Remark2.19. Common-cover condition is proportional a\_i A\_i\textasciicircum{}-1 (Definition2.16), producing Y\_lambda=\allowbreak{}sum y\_j\textasciicircum{}d+\allowbreak{}lambda product y\_j\textasciicircum{}b\_j and generated deck group G. Assume H\textasciicircum{}n(V\_lambda)->H\textasciicircum{}n(V\_lambda\textasciicircum{}star) injective. The exact author characteristic polynomial of Frobenius on H\textasciicircum{}(n-1)(Y\_lambda)\textasciicircum{}G lies in Q[T] and divides each corresponding Xi characteristic polynomial. No unconditional nonsmooth claim or zeta positive-degree consequence is selected.\par
\noindent\textbf{Difficulty:} H2: The proof must descend invariant rigid-cohomology classes through rational monomial quotient maps that are not everywhere morphisms. It constructs auxiliary weighted covers and uses explicit monomial forms and positivity to extend classes across coordinate divisors, then combines compact-support duality with Frobenius-compatible finite-quotient descent. This arithmetic cohomology argument goes beyond citing that varieties share a cover.\par
\noindent\textbf{Editorial boundary note (not author text):} The selected task is only the original Theorem 2.18 final Frobenius characteristic-polynomial rationality/common-factor conclusion, under its retained restriction-map injectivity and the additional source coprimality notice. The full original theorem wording is kept for source fidelity; its separate module quotient-direction assertion is unselected. Complete new rational monomial/torus/weighted auxiliary covers, differential-form bases/reductions, rigid–Monsky–Washnitzer comparison and invariant-form extension across coordinate divisors are bundled as uncounted core. The crucial positivity argument extends forms instead of pretending the rational map is everywhere a morphism. Standard Frobenius rationality, finite-quotient compact-support descent and duality remain cited prerequisites. Original source quotient/subspace wording is unchanged, and no module-direction repair, unconditional nonsmooth claim or separate zeta-function consequence is generated.\par
\noindent\textbf{Source:} \url{https://sigma-journal.com/2017/087/}\quad DOI: \url{https://doi.org/10.3842/SIGMA.2017.087}\par
\noindent\textbf{License and attribution:} Copyright retained by the named authors. Published by SIGMA under CC BY 4.0: \url{https://creativecommons.org/licenses/by/4.0/}. Source: \url{https://sigma-journal.com/about.html}. Adaptation consists of standalone excerpt formatting, cover and numbering restoration; original author language and mathematical text are retained.\par
\noindent\textbf{Editorial symbol notice (not author text):} Introductory1.2 has a matrix/inverse typo; selected2.18 refers to the correct explicit common-cover Definition2.16. Proposition2.15 has a quotient/subspace wording discrepancy; selected statement scope retains only the expressly proved Frobenius characteristic-polynomial conclusion, without adding a module-direction claim.\par
\medskip\hrule\medskip\noindent\textbf{Author text begins below.}\par
\setcounter{section}{1}
\setcounter{subsection}{0}
\setcounter{subsubsection}{0}
\setcounter{Theorem}{2}
\setcounter{equation}{0}
% Original source lines 131--318
\section{Delsarte hypersurface}

Fix an integer $n\geq 2$ and f\/ix a f\/inite f\/ield $\mathbf{F}_q$.
\begin{Definition}An invertible matrix $A:=(a_{i,j})_{0\leq i,j\leq n}$, such that all entries are nonnegative integers is called a \emph{coefficient matrix} if all entries of $A^{-1}(1,\dots,1)^{\rm T}$ are positive and each column of $A$ contains a zero.

In that case let $d$ be an integer such that $B:=dA^{-1}$ has integral coef\/f\/icients. We call $B$ the \emph{map matrix}. We call $B(1,\dots,1)^{\rm T}$ the \emph{weight vector}, which we denote by $\mathbf{w}:=(w_0,\dots,w_n)$.

A vector $\mathbf{a}:=(a_0,\dots,a_n)$ consisting of nonnegative integers such that $\sum\limits_{i=0}^n w_ia_i=d$ holds and such that all entries of $\mathbf{a} A^{-1} $ are nonnegative is called a \emph{deformation vector}.
\end{Definition}

\begin{Definition} Fix a pair $(A,\mathbf{a})$ consisting of coef\/f\/icient matrix and a deformation vector $\mathbf{a}$. Assume that $\gcd(q,d)=1$. Then we call $(A,\mathbf{a})$ \emph{Delsarte deformation data of length~$n$}.
\end{Definition}

Let $(A,\mathbf{a})$ be Delsarte deformation data of length $n$. Let
\begin{gather*} X_\lambda:=Z\left( \sum_{i=0}^n \prod_{j=0}^n x_j^{a_{i,j}} + \lambda \prod_{i=0}^n x_i^{a_i}\right)\end{gather*}
be the corresponding one-parameter family of hypersurfaces of weighted degree $d$ in the weighted projective space $\mathbf{P}(w_0,\dots,w_n)$.

Denote with $(b_0,\dots,b_n)$ the entries of $\mathbf{a} B $. Let $Y_{\lambda}$ be
\begin{gather*} Z\left(\sum_{i=0}^n y_i^d +\lambda \prod_{i=0}^n y_i^{b_i}\right) \subset \mathbf{P}^n.\end{gather*}

\begin{Remark} Our def\/inition of $\mathbf{w}$ may lead to choices of the $w_i$ such that the gcd of $(w_0,\dots,w_n)$ is larger than one. The choice of the $w_i$ is such that the weighted degree of the polynomial def\/ining $X_{\lambda}$ equals the degree of $Y_\lambda$.
\end{Remark}

We have a $(\mathbf{Z}/d\mathbf{Z})^{n}$-action on $\mathbf{P}^n$ induced by
\begin{gather*} (g_1,\dots,g_n)(x_0:x_1:\dots:x_n):=\big(x_0:\zeta^{g_1}x_1:\zeta^{g_2}x_2:\dots:\zeta^{g_n}x_n\big),\end{gather*}
with $\zeta$ a f\/ixed primitive $d$-th root of unity. The subgroup $G$ def\/ined by $\sum\limits_{i=1}^n g_i b_i\equiv 0 \bmod d$ acts on $Y_{\lambda}$.

The rational map $\mathbf{P}^n\dashrightarrow \mathbf{P}(\mathbf{w})$ given by
\begin{gather*} (y_0,y_1,\dots,y_n) \mapsto \left(\prod_{i=0}^n y_i^{b_{0,i}},\prod_{i=0}^n y_i^{b_{1,i}},\dots,\prod_{i=0}^n y_i^{b_{n,i}}\right)\end{gather*}
induces a rational map $Y_{\lambda}\dashrightarrow X_{\lambda}$. This rational map is Galois (i.e., the corresponding extension of function f\/ields is Galois) and the Galois group is a subgroup of~$G$.

In particular, if all the $b_{i,j}$ are nonnegative then this rational map is a morphism. (This map was used by Shioda \cite{ShiExpPic} to give an algorithm to calculate the Picard number of a Delsarte surface in~$\mathbf{P}^3$.)

\begin{Lemma}
The hypersurface $X_0$ is irreducible.
\end{Lemma}
\begin{proof}
 Each column of $A$ contains a zero by the def\/inition of coef\/f\/icient matrix. Hence $x_k$ does not divide \begin{gather*}\sum_{i=0}^n \prod_{j=0}^n x_j^{a_{i,j}}\end{gather*}
for any $k$. Hence for every irreducible component of $X_0$ the points such that all coordinates are nonzero are dense, and these latter points are in the image of $Y_0$. This implies that every irreducible component of $X_0$ is the closure of an irreducible component of the image of $Y_0$. Since $n>1$ it follows that $Y_0$ is irreducible and hence $X_0$ is irreducible.
\end{proof}

\begin{Definition}
We call $X_0$ the \emph{Delsarte hypersurface} associated with $A$ and $X_\lambda$ the \emph{one-dimen\-sio\-nal monomial deformation} associated with $(A,\mathbf{a})$. If, moreover, $X_0$ is quasismooth then we call $X_0$ \emph{invertible hypersurface}.
\end{Definition}

\begin{Example} Consider
\begin{gather*} x_0^4+x_1^4+x_2^3x_3+x_3^3x_2+\lambda x_0x_1x_2x_3.\end{gather*}
Then
\begin{gather*} A=\left(\begin{matrix} 4&0&0&0\\0&4&0&0\\0&0&3&1\\0&0&1&3\end{matrix} \right) \qquad \text{and} \qquad \mathbf{a}=(1,1,1,1).\end{gather*}
We have that
\begin{gather*} B=\left(\begin{matrix} 2&0&0&0\\0&2&0&0\\0&0&3&-1\\0&0&-1&3\end{matrix} \right) \qquad \text{and} \qquad \mathbf{w}=(2,2,2,2). \end{gather*}
In particular, we have that this family is birational a quotient of
\begin{gather*} x_0^8+x_1^8+x_2^8+x_3^8+\lambda (x_0x_1x_2x_3)^2.\end{gather*}
The group $G$ is generated by the automorphisms
\begin{gather*} (x_0,x_1,x_2,x_3)\mapsto(x_0,-x_1,x_2,x_3)\qquad \text{and} \qquad
(x_0,x_1,x_2,x_3)\mapsto\big(x_0,x_1,\zeta^3 x_2,\zeta x_3\big),\end{gather*} with $\zeta$ a primitive $8$-th root of unity.
\end{Example}

\begin{Definition}\label{defGenPos} A hypersurface $X=V(f)\subset \mathbf{P}^n$ is \emph{in general position} if
\begin{gather*} V\left(x_0\frac{\partial}{\partial{x_0}} f, \dots,x_n\frac{\partial}{\partial{x_n}} f\right)\end{gather*}
is empty. Equivalently, $X$ is smooth and for any subset $\{i_1,\dots,i_c\} \subset \{0,1,\dots,n\}$ we have that
\begin{gather*} X \cap V(x_{i_1})\cap \dots \cap V(x_{i_c})\end{gather*}
is also smooth.
\end{Definition}

\begin{Lemma}\label{lemGenPos} If $Y_\lambda$ is smooth then $Y_{\lambda}$ is in general position.
\end{Lemma}

\begin{proof} Suppose we intersect $Y_{\lambda}$ with $x_{i_1}=\dots=x_{i_c}=0$. If some $b_{i_j}$ is nonzero then the inter\-section is a Fermat hypersurface in $\mathbf{P}^{n-c}$ and is smooth. If all $b_{i_j}$ are zero then we can do the following: After a change of coordinates we may assume that $\{i_1,\dots,i_{c}\}=\{0,1,\dots,c-1\}$. We now have that $Y_{\lambda}$ is the zero set of
\begin{gather*} \sum_{i=0}^{c-1} x_i^d+h(x_c,\dots,x_n)\end{gather*}
for some $h\in \mathbf{F}_q[x_c,\dots,x_n]$. From $\gcd(q,d)=1$ it follows that the singular points of the intersection $V(x_0,\dots,x_{c-1},h)$ are in one-to-one correspondence with the singular points of $Y_\lambda$. Hence $V(x_0,\dots,x_{c-1},h)$ is smooth.
\end{proof}

Recall that we started with a hypersurface $X_\lambda \subset \mathbf{P}(\mathbf{w})$ and constructed a hypersurface $Y_{\lambda}\subset \mathbf{P}^n$, such that $X_\lambda$ is birational to a quotient of $Y_\lambda$. Denote with $U_{\lambda}:=\mathbf{P}(\mathbf{w})\setminus X_\lambda$ and $V_{\lambda}:=\mathbf{P}^n\setminus Y_\lambda$ be the respective complements.

Denote now with $(\mathbf{P}(\mathbf{w}))^*$, $U_{\lambda}^*$, $V_{\lambda}^*$, $X_{\lambda}^*$, $Y_\lambda^*$, etc. the original variety minus the intersection with $Z(x_0\dots x_n)$ or $Z(y_0 \dots y_n)$, the union of the coordinate hyperplanes. We have that the quotient map $\mathbf{P}^n \dashrightarrow \mathbf{P}(\mathbf{w})$ def\/ines surjective morphisms $(\mathbf{P}^n)^*\to \mathbf{P}(\mathbf{w})^*$, $Y_\lambda^*\to X_{\lambda}^*$, $V_\lambda^*\to U_\lambda^*$.

There is a second quotient map $\mathbf{P}^n \to \mathbf{P}(\mathbf{w})$ given by $(z_0:\dots:z_n) \to (z_0^{w_0}:\dots:z_n^{w_n})$. This map is a morphism and is a ramif\/ied Galois covering. Denote with $H$ the corresponding Galois group. Let $\tilde{X}_{\lambda}$ be the pull back of $X_\lambda$ and let $\tilde{U}_{\lambda}$ be the pull back of $U_{\lambda}$.

Fix now a lift $\mu\in \mathbf{Q}_q$ of $\lambda$. Then we can def\/ine $F_{\mu}$, $\tilde{U}_{\mu}$, $V_{\mu}$, $\tilde{ X}_{\mu}$, $Y_{\mu}$ similarly as above. If $y_0,\dots,y_n$ are projective coordinates on $\mathbf{P}^n$ then let $\Omega$ be
\begin{gather*} \left(\prod_{i=0}^n y_i\right) \left(\sum_{i=0}^n (-1)^{i} \frac{dy_0}{y_0} \wedge \dots \wedge \widehat{\frac{dy_i}{y_i} }\wedge\dots \wedge\frac{dy_n}{y_n} \right).\end{gather*}
We recall now some standard notation used to study the cohomology of a hypersurface complement in $\mathbf{P}^n$.

\begin{notation}Let $\mathbf{m}=(m_0,\dots,m_n)$ be $(n+1)$-tuple of positive integers, such that $\sum\limits_{i=0}^n m_i =td$ for some positive integer~$t$. Then
\begin{gather*} \tilde{\omega}_{\mathbf{m}}:= \frac{\prod\limits_{i=0}^n x_i^{m_i-1}}{(F_{A,\mu}(x_0^{w_0},\dots,x_n^{w_n}))^t} \Omega\end{gather*}
is an $n$-form on the complement $\tilde{U}_{\mu}$ of $\tilde{X}_{\mu}$. If we allow the $m_i$ and $t$ to be arbitrary integers such that the equality $\sum\limits_{i=0}^n m_i=td$ holds then $\tilde{\omega}_{\mathbf{m}}$ is a form on $\tilde{U}^*_{\mu}$.

Let $\mathbf{m}=(m_0,\dots,m_n)$ be $(n+1)$-tuple of positive integers, such that $\sum\limits_{i=0}^n m_i =td$ holds for some positive integer $t$. Let $D$ be the diagonal matrix $d I_{n+1}$. Then
\begin{gather*} \omega_{\mathbf{m}}:= \frac{\prod\limits_{i=0}^n y_i^{m_i-1}}{F_{D,\mu}^t} \Omega\end{gather*}
is an $n$-form on the complement $V_{\mu}$ of $Y_{\mu}$. If we allow the $m_i$ and $t$ to be arbitrary integers such that the equality $\sum\limits_{i=0}^n m_i=td$ then $\omega_{\mathbf{m}}$ is a form on~$V^*_{\mu}$.
\end{notation}

The following result seems to be known to the experts, but we include it for the reader's convenience:

\begin{Lemma}\label{lemBasis} There exists a finite set $S\subset \mathbf{Q}_q$ such that $0\not \in S$ and for all $\mu\in \mathbf{Q}_q\setminus S$ we have that
\begin{gather*}\mathcal{B}:=\left \{\omega_{\mathbf{m}} \colon 0<m_i<d \mbox{ for } i=0,\dots, n \ \text{and} \ \sum_{i=0}^n m_i\equiv 0 \bmod d\right \}\end{gather*}
is a basis for $H^n_{\dR}(V_\mu,\mathbf{Q}_q)$.

Similarly, there exists a finite set $S^*$ such that $0\not \in S^*$ and for all $\mu\in\mathbf{Q}_q\setminus S^*$ we have that
\begin{gather*}\mathcal{B}^*:=\left \{\omega_{\mathbf{m}} \colon 0\leq m_i<d \mbox{ for } i=0,\dots, n \ \text{and} \ \sum_{i=0}^n m_i\equiv 0 \bmod d\right\}\end{gather*}
is a basis for $H^n_{\dR}(V_\mu^*,\mathbf{Q}_q)$.
\end{Lemma}
\begin{proof} The forms $\omega_{\mathbf{m}}$, such that $m_i \geq 1$ for $i=0,\dots,n$ generate the de Rham cohomology group $H^n_{\dR}(V_\mu)$.
By dif\/ferentiating certain particular $(n-1)$-forms on $V_\mu$ we have that the following relation in $H^n_{\dR}(V_\mu)$
\begin{gather}\label{eqnRed} \frac{G_{y_i}}{F^t} \Omega = \frac{t G F_{y_i}}{F^{t+1}} \Omega \end{gather}
for any form $G\in \mathbf{Q}_q[y_0,\dots,y_n]_{td-n}$. (This is the so-called Grif\/f\/iths--Dwork method to reduce forms in cohomology.)

For $\mu=0$ we have that $F_{y_i}=dy_i^{d-1}$. Using (\ref{eqnRed}) we f\/ind the relation \begin{gather} \label{eqnRedb} \frac{y_0^{m_0} G(y_1,\dots,y_n)}{F_0^{t+1}}\Omega = \frac{(m_0-d+1)y_0^{m_0-d} G(y_1,\dots,y_n)}{t F_0^t} \Omega\end{gather}
and similar relations for the other $y_i$. In this way we can reduce forms such that all exponents are at least $0$ and at most $d-1$. However, if an exponent equals $d-1$ then this relation yields that the class is zero in cohomology. In particular, the $\omega_{\mathbf{m}}$ with $0<m_i<d$ for $i=0,\dots, n$ and $\sum\limits_{i=0}^n m_i\equiv 0 \bmod d$ generate $H^n_{\dR}(V_0)$. Grif\/f\/iths \cite{GriRat} showed that the relations of type~(\ref{eqnRed}) generate all relations and hence $\mathcal{B}$ is a basis for $H^n_{\dR}(V_0)$. If $X_\mu$ is smooth then the dimension of $H^n_{\dR}(V_\mu)$ is independent of $\mu$ and it is then straightforward to check that there are at most f\/initely many choices of $\mu$ for which $X_\mu$ is smooth and $\mathcal{B}$ is not a basis for $H^n_{\dR}(V_\mu)$.

We now prove the statement on $H^n_{\dR}(V_\mu^*)$. Note that if $X_\mu$ is smooth then by Lemma~\ref{lemGenPos} it is in general position. Therefore the dimension of $H^n_{\dR}(V_\mu^*)$ is independent of $\mu$. Hence it suf\/f\/ices to show that $\mathcal{B}^*$ is a basis for $H^n_{\dR}(V_0^*)$. Again we have relations of type~(\ref{eqnRed}), but now we may take $G\in \mathbf{Q}_q\big[y_0,y_0^{-1},\dots,y_n,y_n^{-1}\big]_{td-n}$. If the exponent of a~variable is at most $-2$ then we can use the relations of the shape~(\ref{eqnRedb}) to increase the exponent of this variable. However, if the exponent equals $-1$ we cannot do this, because then we would have to divide by zero in~(\ref{eqnRedb}). In this way we obtain that $\mathcal{B}^*$ generates $H^n_{\dR}(V^*_0)$. Moreover, as in the above case there are no further relations and~$\mathcal{B}^*$ is a basis.
\end{proof}

\begin{Remark} The $\mu$ for which $\mathcal{B}$ is not a basis can be determined by the methods of \cite[Section~3]{PanTui}.
\end{Remark}

\begin{Remark}Denote with $H^n_{\MW}(V_{\lambda})$ the $n$-th Monsky--Washnitzer cohomology of $V_\lambda$. The Monsky--Washnitzer cohomology is essentially the cohomology of the tensor product of de~Rham complex of a lift of $V_{\lambda}$ to characteristic zero with a weakly complete f\/initely generated algebra~$A^{\dagger}$. The Frobenius action on the cohomology is induced by a lift of Frobenius to $A^{\dagger}$. For more details see \cite[Theorem~2.4.5]{PutMW}. In that paper it is shown that two dif\/ferent lifts of $V_{\lambda}$ yield isomorphic complexes and two choices of lifts of Frobenius yield homotopic maps on the complexes. In particular, $H^n_{\MW}(V_\lambda)$ is independent of the choices made.

Let $\mu\in \mathbf{Q}_q$ be a lift of $\lambda$. One choice of a lift of $V_\lambda$ to characteristic zero is $V_{\mu}$, and the construction of the Monsky--Washnitzer cohomology yields a natural map $H^n_{\dR}(V_\mu)\to H^n_{\MW}(V_\lambda)$ of $\mathbf{Q}_q$-vector spaces. If $Y_\lambda$ is smooth then the is an isomorphism by \cite{BalChi}. Since $V_\lambda \subset \mathbf{P}^n$ is af\/f\/ine and smooth we have an isomorphism $H^n_{\MW}(V_\lambda)\cong H^n_{\rig}(V_\lambda)$ by~\cite{BerSMF}, where the latter group is rigid cohomology. Since there are inf\/initely many lifts~$\mu$ of $\lambda$ we can always choose a lift~$\mu$ such that~$\mathcal{B}$ is a basis for $H^n_{\dR}(V_\mu)$ and thereby yielding a basis for $H^n_{\MW}(V_{\lambda})$.

If $Y_\lambda$ is smooth then using Lemma~\ref{lemGenPos} we f\/ind that $V_\lambda^*$ is the complement of a normal crossing divisor. In particular, we can apply \cite{BalChi} and f\/ind a natural isomorphism $H^n_{\MW}(V_\lambda^*)\cong H^n_{\dR}(V_\mu^*)$. As above, we have an isomorphism $H^n_{\MW}(V_\lambda^*)\cong H^n_{\rig}(V_\lambda^*)$ and we identif\/ied a basis for $H^n_{\rig}(V_\lambda^*)$.
\end{Remark}

\begin{Remark} The action of $G$ lifts to characteristic zero. The forms $\omega_{\mathbf{m}}$ are eigenvectors for~$g^*$ each element $g\in G$. Hence the $G$-invariant ones span $H^n_{\MW}(V_\lambda)^G$.

If $Y_{\lambda}$ is singular then by the def\/inition of Monsky--Washnitzer cohomology we have that $H^n_{\MW}(V_\lambda)$ is generated by expressions
\begin{gather*} \sum_{\mathbf{m}=(m_0,\dots,m_n),m_i\geq 1} a_{\mathbf{m}} \omega_{\mathbf{m}}\end{gather*}
such that there exists $c_1,c_2 \in \mathbf{Q}$ with $c_1>0$ and $v(a_{\mathbf{m}})\geq c_1\big(\sum\limits_{i=0}^n m_i\big)+c_2$. The space $H^n_{\MW}(V_\lambda^*)$ is generated by expression
\begin{gather*} \sum_{\mathbf{m}=(m_0,\dots,m_n),m_i\geq -N} a_{\mathbf{m}} \omega_{\mathbf{m}}\end{gather*}
such that there exists $c_1,c_2 \in \mathbf{Q}$ with $c_1>0$ and $v(a_{\mathbf{m}})\geq c_1\big(\sum\limits_{i=0}^n m_i\big)+c_2$. The $G$-invariant subspaces are generated by similar sums, but in which only the $G$-invariant~$\omega_{\mathbf{m}}$ occur.
\end{Remark}

\begin{Remark} If one wants to study the Frobenius matrix by using the dif\/ferential equations, like in~\cite{Katz} or in~\cite{DKSSVW} then one needs to be more careful in lifting $V_\lambda$ to characteristic zero. In~\cite{Katz} one has to take~$\mu$ to be the Teichm\"uller lift of~$\lambda$. The reason for this, is that a priori Frobenius maps $H^i\big(U_\mu^q\big)$ to $H^i(U_\mu)$. To have an operator on $H^n(U_\mu)$ we need that $\mu^q=\mu$. If one works directly with Frobenius on Monsky--Washnitzer chomology then this constraint on $\mu$ does not exist.
\end{Remark}

From now on we use $H^i$ and $H^i_c$ to indicate rigid cohomology respectively rigid cohomology with compact support.

By \cite[Proposition 2.1]{Zetafamrec} we have canonical isomorphisms
\begin{gather*} H^i_c(U_\lambda^*)\cong H^i_c(V_\lambda^*)^G \qquad \text{and} \qquad H^i_c(U_\lambda^*)\cong H^i_c\big(\tilde{U}_\lambda^*\big)^H.\end{gather*}
We want to compare the cohomology of $H^n_c(V_\lambda)^G$ with the cohomology of~$H^n_c(U_{\lambda})$. However, $U_\lambda$ may be singular, hence we work with $H^n_c\big(\tilde{U}_\lambda\big)^ H$ instead. Using Poin\-car\'e duality it suf\/f\/ices to compare $H^n(V_\lambda)^ G$ with $H^n\big(\tilde{U}_\lambda\big)^H$ instead. Since both varieties are smooth and af\/f\/ine we can identify their rigid cohomology groups with their Monsky--Washnitzer cohomology groups. We will do this in order to prove:

\begin{Proposition}\label{prpFactor} Suppose that $H^n(V_\lambda)\to H^n(V_{\lambda}^*)$ is injective. Then $H^n(V_{\lambda})^G$ is a quotient of $H^n\big(\tilde{U}_{\lambda}\big)^H$. In particular, the characteristic polynomial of Frobenius acting on $H^n_c(V_{\lambda})^G$ is in~$\mathbf{Q}[T]$ and divides the characteristic polynomial of Frobenius acting on~$H^n_c(U_{\lambda})$.
\end{Proposition}

\begin{proof}Since $H^n(V_\lambda)\to H^n(V_{\lambda}^*)$ is injective we have by \cite{BerPoi} that the Poincar\'e dual of this map is surjective, and therefore that $H^n_c(V_\lambda)^G$ is a quotient of~$H^n_c(V_\lambda^*)^G$. This implies that $H^n_c(V_\lambda)^G$ is also a quotient of $H^n_c\big(\tilde{U}^*_\lambda\big)^H$. Hence it suf\/f\/ices to show that the kernel of natural map $H^n_c\big(\tilde{U}^*_\lambda\big)\to H^n_c\big(\tilde{U}_\lambda\big)$ is mapped to zero in $H^n_c(V_\lambda)^G$. Using Poincar\'e duality we can consider $H^n(V_\lambda)^G$ as a subspace of $H^n\big(\tilde{U}_{\lambda}^*\big)^H$. It suf\/f\/ices to show that $H^n(V_\lambda)^G$ is in the image of~$H^n(\tilde{U}_\lambda)$.

A form $\omega_{\mathbf{k}}$ is in $H^n(V_\lambda)^G$ if and only if there is a monomial type $\mathbf{m}_0$ such that $\mathbf{k}=\mathbf{m}_0 B$. We identif\/ied $H^n(V_\lambda)^G$ with a subspace of $H^n\big(\tilde{U}_\lambda^*\big)^H$. The class of $\omega_{\mathbf{k}}$ is identif\/ied with $\tilde{\omega}_{\mathbf{m}}$ where $\mathbf{m}=\mathbf{m}_0(\diag(w_0,\dots,w_n))$.

The entries of $\mathbf{m}$ are integers, which may be nonpositive. If all entries of $\mathbf{m}$ are positive then~$\omega_{\mathbf{m}}$ is in the image of $H^n\big(\tilde{U}_{\lambda}\big)^ H$. Recall that $B=d A^{-1}$ and therefore $\mathbf{m}_0=\mathbf{k} \frac{1}{d} A $. Since~$\mathbf{k}$ has positive entries, $A$ has positive entries and no zero column it follows that also the entries of~$\mathbf{m}_0$ are positive and therefore all entries of~$\mathbf{m}$ are also positive. This yields the f\/irst statement.

To prove the second statement. By \cite[Lemma~4.3]{Zetafamrec} it follows that $H^n_c(V_\lambda)^G$ is Frobenius invariant and the characteristic polynomial is in $\mathbf{Q}[T]$. Using Poincar\'e duality we f\/ind that~$H^n_c(V_\lambda)^G$ is a subspace of~$H^n_c(\tilde{U}_\lambda)^H$. As explained above, the latter space is isomorphic with~$H^n_c(U_\lambda)$.
\end{proof}

\begin{Definition}
Fix Delsarte deformation data $(A_1,\mathbf{a}_1), \dots, (A_t,\mathbf{a}_t)$ of length~$n$. We say that they have a \emph{common cover}
if for every $i$, $j$ we have that $\mathbf{a}^{\rm T}_i A_i^{-1}$ and $\mathbf{a}^{\rm T}_j A_j^{-1}$ are proportional.
\end{Definition}
\begin{Example} Take the following f\/ive matrices
\begin{gather*}
\left(\begin{matrix} 4&0&0&0\\0&4&0&0\\0&0&4&0\\0&0&0&4\end{matrix} \right)\!,\quad
\left(\begin{matrix} 4&0&0&0\\0&4&0&0\\0&0&3&1\\0&0&1&3\end{matrix} \right)\!,\quad
\left(\begin{matrix} 3&1&0&0\\1&3&0&0\\0&0&3&1\\0&0&1&3\end{matrix} \right)\!,\quad
\left(\begin{matrix} 4&0&0&0\\0&3&1&0\\0&0&3&1\\0&1&0&3\end{matrix} \right)\!,\quad
\left(\begin{matrix} 3&1&0&0\\0&3&1&0\\0&0&3&1\\1&0&0&3\end{matrix} \right)\!.
\end{gather*}
In each case we take $(1,1,1,1)^{\rm T}$ as the deformation vector then the $(A_i,\mathbf{a}_i)$ have a common cover.
\end{Example}

Suppose now that $(A_1,\mathbf{a}_1),\dots,(A_t,\mathbf{a}_t)$ have a common cover. Let $d$ be the smallest positive integer such that $dA_i^{-1}$ has integral coef\/f\/icients for all~$i$. The sum of the entries of $\mathbf{b}_i=\mathbf{a}_i \big(dA_i^{-1}\big)$ equals~$d$. By assumption we have that for each $i$ and $j$ the vectors $\mathbf{b}_i$ and $\mathbf{b}_j$ are proportional, hence these vectors coincide and we denote this common vector with~$\mathbf{b}$.

Denote with $b_j$ the entries of $\mathbf{b}$. Denote with $X_{i,\lambda}$ the family associated with $(A_i,\mathbf{a}_i)$. Then $X_{i,\lambda}$ is birational to a quotient of
\begin{gather*}Y_\lambda\colon \ \sum_{i=0}^n y_i^{d} +\lambda \prod_{i=0}^n y_i^{b_i}.\end{gather*}
At the beginning of this section we gave an explicit description of this map. From that description it follows that $Y_\lambda \to X_{i,\lambda}$ is def\/ined whenever all the~$y_i$ are nonzero.

We can now apply Proposition~\ref{prpFactor} to the above setup and we f\/ind directly that:

\begin{Theorem} \label{thmMainO} Let $(A_1,\mathbf{a}_1), \dots, (A_t,\mathbf{a}_t)$ be Delsarte deformation data of length $n$ with a common cover. Denote with~$X_{i,\lambda}$ be the corresponding families of Delsarte hypersurfaces and with~$Y_\lambda$ the common cover. Let $G_i$ be the Galois group of the function field extension corresponding to the rational map $Y_\lambda\dashrightarrow X_{i,\lambda}$. Then the automorphisms in~$G_i$ extend to automorphisms of~$Y_{\lambda}$. Identify~$G_i$ with the corresponding subgroup of $\Aut(Y_{\lambda})$. Let $G=G_1.G_2.\dots.G_t\subset \Aut(Y_\lambda)$.

Suppose that $H^n(V_\lambda)\to H^n(V_\lambda^*)$ is injective. Then for each $i=1,\dots,t$ we have that $H^n_c(V_{\lambda})^G$ is a quotient of $H^n_c(U_{i,\lambda})$. In particular, the characteristic polynomial of Frobenius on~$H^{n-1}(Y_\lambda)^G$ is in $\mathbf{Q}[T]$ and is a common factor of the characteristic polynomials of Frobenius acting on~$H^{n-1}(X_{i,\lambda})$.
\end{Theorem}

\begin{Remark} Recall that in order to be Delsarte deformation data we need that $\gcd(q,(n+1)\det(A_i))=1$ for all $i$.
\end{Remark}

\begin{Remark} If $Y_\lambda$ is smooth then it is in general position by Lemma~\ref{lemGenPos}.

The map $H^ {n-1}(Y_{\lambda})\to H^{n-1}(Y_{\lambda}^*)$ is injective if $n-1$ is even, and has a kernel if $n-1$ is odd, and this kernel is generated by the hyperplane class, see \cite[Theorem~1.19]{Katz}. The residue map identif\/ies $H^{n}(V_{\lambda})$ with the primitive part of the cohomology of $H^{n-1}(Y_{\lambda})$. In particular, the composition $H^n(V_{\lambda})\to H^{n-1}(Y_{\lambda}^*)$ is injective independent of the parity of $n$. From the diagram on \cite[p.~79]{Katz} it follows that the latter map factors through $H^n(V_\lambda^*)$. In particular, $H^n(V_\lambda)\to H^n(V_\lambda^*)$ is injective. Hence we can apply the above proposition if $Y_\lambda$ is smooth. The values of $\lambda$ for which $Y_\lambda$ is singular can be determined from the formula \cite[Lemma~3.7]{zetafam}.
\end{Remark}
\begin{thebibliography}{99}
\bibitem[1]{BalChi}
Baldassarri F., Chiarellotto B., Algebraic versus rigid cohomology with
 logarithmic coef\/f\/icients, in Barsotti {S}ymposium in {A}lgebraic {G}eometry
 ({A}bano {T}erme, 1991), \textit{Perspect. Math.}, Vol.~15, Academic Press,
 San Diego, CA, 1994, 11--50.

\bibitem[2]{BerSMF}
Berthelot P., G\'eom\'etrie rigide et cohomologie des vari\'et\'es
 alg\'ebriques de caract\'eristique {$p$}, in Study Group on Ultrametric
 Analysis, 9th year: 1981/82, {N}o. 3 ({M}arseille, 1982), Inst. Henri
 Poincar\'e, Paris, 1983, Exp. No.~J2, 18~pages.

\bibitem[3]{BerPoi}
Berthelot P., Dualit\'e de {P}oincar\'e et formule de {K}\"unneth en
 cohomologie rigide, \href{https://doi.org/10.1016/S0764-4442(97)88895-7}{\textit{C.~R.~Acad. Sci. Paris S\'er.~I Math.}}
 \textbf{325} (1997), 493--498.

\bibitem[8]{DKSSVW}
Doran C.F., Kelly T.L., Salerno A., Sperber S., Voight J., Whitcher U., Zeta
 functions of alternate mirror {C}alabi--{Y}au families, \href{https://arxiv.org/abs/1612.09249}{arXiv:1612.09249}.

\bibitem[11]{GriRat}
Grif\/f\/iths P.A., On the periods of certain rational integrals.~{II},
 \href{https://doi.org/10.2307/1970747}{\textit{Ann. of Math.}} \textbf{90} (1969), 496--541.

\bibitem[12]{Katz}
Katz N.M., On the dif\/ferential equations satisf\/ied by period matrices,
 \href{https://doi.org/10.1007/BF02698924}{\textit{Inst. Hautes \'Etudes Sci. Publ. Math.}} (1968), 223--258.

\bibitem[15]{zetafam}
Kloosterman R., The zeta function of monomial deformations of {F}ermat
 hypersurfaces, \href{https://doi.org/10.2140/ant.2007.1.421}{\textit{Algebra Number Theory}} \textbf{1} (2007), 421--450,
 \href{https://arxiv.org/abs/math.NT/0703120}{math.NT/0703120}.

\bibitem[16]{Zetafamrec}
Kloosterman R., Group actions on rigid cohomology with compact support --
 {E}rratum to ``The zeta function of monomial deformations of {F}ermat
 hypersurfaces'', {P}reprint.

\bibitem[18]{PanTui}
Pancratz S., Tuitman J., Improvements to the deformation method for counting
 points on smooth projective hypersurfaces, \href{https://doi.org/10.1007/s10208-014-9242-8}{\textit{Found. Comput. Math.}}
 \textbf{15} (2015), 1413--1464, \href{https://arxiv.org/abs/1307.1250}{arXiv:1307.1250}.

\bibitem[20]{ShiExpPic}
Shioda T., An explicit algorithm for computing the {P}icard number of certain
 algebraic surfaces, \href{https://doi.org/10.2307/2374678}{\textit{Amer.~J. Math.}} \textbf{108} (1986), 415--432.

\bibitem[21]{PutMW}
van~der Put M., The cohomology of {M}onsky and {W}ashnitzer, \href{https://doi.org/10.24033/msmf.324}{\textit{M\'em.
 Soc. Math. France (N.S.)}} (1986), 33--59.

\end{thebibliography}
\end{document}
